\begin{exercise}[Uniqueness of the Levi--Civita connection]{Tong, Cambridge Part III Sheet 2}
Starting only from zero torsion and metric compatibility, derive the
Christoffel formula for \(\Gamma^\rho{}_{\mu\nu}\).  Prove uniqueness by
considering the difference of two such connections.
\end{exercise}

\begin{exercise}[Covariant constancy of the volume form]{Tong, Cambridge Part III Sheet 2}
For the Levi--Civita tensor
\(\varepsilon_{\mu_1\ldots\mu_n}=\sqrt{|g|}\,[\mu_1\ldots\mu_n]\), prove
\(\nabla_\rho\varepsilon_{\mu_1\ldots\mu_n}=0\).  Use this to show that the
Hodge dual commutes with the covariant derivative.
\end{exercise}

\begin{exercise}[The curved-space divergence theorem]{Tong, Cambridge Part III Sheet 1}
Let \(V^\mu\) be a vector on a compact region \(\mathcal U\) with boundary
\(\partial\mathcal U\).  Derive
\[
\int_{\mathcal U}\sqrt{|g|}\,\nabla_\mu V^\mu\,\dd^nx
=\int_{\partial\mathcal U}V^\mu n_\mu\,\dd\Sigma
\]
and state the orientation conventions needed in Lorentzian signature.
\end{exercise}

\begin{exercise}[Frobenius and hypersurface orthogonality]{Cambridge Part III, General Relativity examples}
Let \(u_\mu\) be a nowhere-vanishing one-form.  Prove locally that
\(u_\mu=f\,\partial_\mu t\) for some \(f,t\) if and only if
\(u\wedge\dd u=0\).  Express the condition in index notation and relate it to
the vorticity of a timelike congruence.
\end{exercise}

\begin{exercise}[Cartan equations in polar coordinates]{MIT 8.962, Problem Set 6}
On the Euclidean plane use the orthonormal coframe
\(e^1=\dd r\), \(e^2=r\,\dd\phi\).  Solve the torsion-free Cartan equation for
the connection one-form and use the second Cartan equation to verify that the
curvature vanishes despite the nonzero connection.
\end{exercise}

\begin{exercise}[Connection under a conformal rescaling]{Cambridge Part III, 2025 exam}
For \(\widetilde g_{\mu\nu}=e^{2\omega}g_{\mu\nu}\), derive
\(\widetilde\Gamma^\rho{}_{\mu\nu}-\Gamma^\rho{}_{\mu\nu}\).  Show directly
that null geodesics are unchanged as unparameterized curves.
\end{exercise}

\begin{exercise}[Metric expansion in normal coordinates]{MIT 8.962, Problem Set 7}
In Riemann normal coordinates centered at \(p\), show that
\(\Gamma^\rho{}_{\mu\nu}(p)=0\) and derive the quadratic metric expansion
\[
g_{\mu\nu}(x)=\eta_{\mu\nu}
-\frac13R_{\mu\alpha\nu\beta}(p)x^\alpha x^\beta+\order(x^3).
\]
Explain why curvature cannot be transformed away.
\end{exercise}

\begin{exercise}[Lie brackets from noncommuting flows]{Tong, Cambridge Part III Sheet 1}
Let \(\Phi_s\) and \(\Psi_t\) be the local flows generated by vector fields
\(X\) and \(Y\).  Follow a point around the infinitesimal loop
\(\Phi_{-s}\Psi_{-t}\Phi_s\Psi_t\) and show that the leading displacement is
proportional to \(st[X,Y]\).  Interpret the result geometrically.
\end{exercise}
